\section{Population changes are new analyses}

Two consumers can agree on an action and outcome while referring to different
customers. A marketing sample might overrepresent a low-risk stratum, while a
servicing portfolio contains mostly a high-risk stratum. Changing the
population field in a copied result would not transport the effect. The
inspected FinMarketing comparison requires equality of population, action,
comparator, origin, horizon, finance owner and capital convention when mapping
an estimate into value; its comment explicitly distinguishes a unit conversion
from reaggregation across populations or horizons~\citep{exec13_i3_inspection}.
This gives the proposed adapter a useful boundary: a new population needs
a new target identity and a scientific bridge.

For a simple bridge, let $X$ have two strata, let $P$ describe the source
population and $Q$ the target population, and suppose the stratum-specific
causal effect $\tau(x)$ is invariant between them. This is an explicit
transport assumption in addition to identification of the source effects.
If $Q(x)>0$ only where $P(x)>0$, then
\begin{equation}
 \tau_Q=\sum_xQ(x)\tau(x)
       =\sum_xP(x)\frac{Q(x)}{P(x)}\tau(x).
 \label{exec13:i3-transport}
\end{equation}
The equality is algebraic under the stated assumptions. It does not establish
that equal stratum names represent equal measurements, that treatment versions
match, or that effects are in fact invariant. Different credit limits,
collection actions or outcome definitions would require a richer bridge.

Suppose the stratum effects are $10$ and $-10$ USD per account. With source
weights $(0.8,0.2)$ the source average is $6$ USD. With target weights
$(0.2,0.8)$ the target average is $-6$ USD. Density ratios $(1/4,4)$
reproduce the second value from the source distribution exactly~\citep{exec13_i3_composition}. The sign reversal arises entirely from the
population change, despite identical conditional effects. A common API and
identical units would not prevent it. If the second stratum is absent from
the source, the target result is not obtained by capping the missing ratio.
One must restrict the target, supply additional information or retain a bound
under stated restrictions.

For a FinMarketing or SMEwallet consumer, a proposed export should therefore
identify the unit of analysis, stratum construction, population counts or
weights, intervention versions and availability of the variables used to
transport. If estimated effects and estimated population weights are combined,
their joint uncertainty matters. Splitting the two estimates across
repositories does not make their errors independent. A consumer can instead
request a conditional planning calculation with fixed supplied weights, but
the response should preserve that qualification. The public-data and transport
chapters discuss alternatives when the desired bank population is not yet
observed; a bank-specific transported claim is not needed to develop this
proposal.

\section{Portfolio composition and the ownership of cash flows}

Separate funding and servicing effects also need not add. Consider a
generated net-value table for one common account, with funding action $F$
and servicing action $S$:
\[
 Y(f,s)=100+20f+15s-30fs,\qquad f,s\in\{0,1\}.
\]
All entries are stipulated economic net values at one common horizon. The
funding-only effect is $20$, the servicing-only effect is $15$, and their
joint effect is $5$. In general the decomposition is
\begin{align}
 E\{Y(1,1)-Y(0,0)\}
   &=E\{Y(1,0)-Y(0,0)\}+E\{Y(0,1)-Y(0,0)\}+I,\\
 I&=E\{Y(1,1)-Y(1,0)-Y(0,1)+Y(0,0)\}.
 \label{exec13:i3-interaction}
\end{align}
Here $I=-30$. Adding two correct isolated analyses and reporting $35$
would omit the interaction. The example does not assert that this particular
interaction occurs in a bank; it demonstrates information the proposed
portfolio consumer needs before combining answers.

Suppose the funding population has 100 accounts, the servicing population
60, and 20 accounts receive both actions. Assume the table applies to each
account, memberships are fixed before treatment, and accounts do not affect
one another. The total incremental value is
\[
 80(20)+40(15)+20(5)=2300,
\]
equivalently $100(20)+60(15)+20(-30)$. The naive separate total is
$2900$. Both evaluations and the equivalent decomposition are reproduced
exactly~\citep{exec13_i3_composition}. Origin identities are consequently
needed for more than deduplication: they establish which units receive the
joint policy. If linkage is uncertain, both membership and the relevant
interaction can be uncertain. If there are spillovers, the no-interference
assumption used by this sum must be replaced rather than hidden in the
aggregation code.

Before adding monetary outputs, the consumer must also identify whose cash
flow is being counted. An internal transfer of $3$ USD recorded as income by
a branch and cost by treasury sums to zero for the consolidated bank. The
branch target and bank target are different even though both are expressed in
USD. Similarly, a principal repayment may be relevant to liquidity while not
being income under the declared net-contribution measure. A state-space
quantity called a balance and a financial quantity called contribution cannot
be combined merely because both have currency units.

The inspected FinMarketing \texttt{CashEntry} distinguishes revenue, cost,
principal, liability, internal transfer and continuation, and retains owner,
cash date, economic basis and inclusion in value. Its \texttt{Value}
validation restricts owner and currency to the target, disallows principal,
liability and internal transfer from the net-contribution measure, and
requires continuation to start after the explicit cash horizon~\citep{exec13_i3_composition}. These are concrete source declarations for a
proposed valuation bridge. They do not establish all accounting policies or
cross-consumer correctness. CardNPV's continuation convention still requires
its own mapping, and a loan cash-flow valuation may legitimately include
principal under a different explicitly declared target.

A proposed combined workflow should therefore retain event and account
identity, economic owner, category, horizon and intervention membership before
aggregation. A manual reference can reconcile one small joint scenario across
the consumer records. An automated adapter should reproduce that result and
retain the same exceptions when an owner, horizon or intervention changes.
When a bridge cannot be justified, separate displayed values with a clear
non-additivity qualification are more informative than a fabricated total.
The resulting next experiment is specific: determine whether the consequential
joint response or population bridge is adequately supported, rather than
repeating isolated analyses whose individual accuracy cannot resolve it.
